%% ch4_method.tex  ?? Chapter 4: The HSP-Agile Pipeline: Design and Theory
%% HSP-Agile Thesis, East China Normal University, 2026
%% Ye Xujun

\chapter{The HSP-Agile Pipeline: Design and Theory}
\label{ch:method}

This chapter gives a complete technical account of the HSP-Agile pipeline.
We first articulate the five design principles that shaped every architectural
decision (Section~\ref{sec:method:principles}), then present the pipeline at
three levels of abstraction: a high-level data-flow view
(Section~\ref{sec:method:overview}), followed by detailed treatments of each
stage: LLM-based candidate synthesis (Section~\ref{sec:method:synthesis}),
SMT-driven formal conformance checking (Section~\ref{sec:method:checker}),
counterexample-guided repair (Section~\ref{sec:method:repair}), and the
requirement pattern guard (Section~\ref{sec:method:guard}).
Section~\ref{sec:method:accept} introduces the hybrid acceptance predicate
that combines the outputs of all four stages into a single pass/fail decision.
The chapter closes with formal proofs of termination and soundness
(Section~\ref{sec:method:theory}).

Throughout this chapter we inherit the notation introduced in
Chapter~\ref{ch:formalization}: $t = (\mathcal{S}_t, \mathcal{X}_t,
\mathcal{Y}_t)$ is an Agile-SOFL task, $\mathcal{S}_t = \langle
s_1,\ldots,s_n\rangle$ is its ordered FSF specification, and $g_t :
\mathcal{X}_t \to \mathcal{Y}_t$ is the specification oracle defined by
first-match semantics.

% ---------------------------------------------------------------------------
\section{Design Principles}
\label{sec:method:principles}
% ---------------------------------------------------------------------------

The design of HSP-Agile is motivated by five principles, each addressing a
concrete limitation of existing automated software engineering approaches
\citep{woodcock2009formalsurvey, chen2021codex}.

\paragraph{P1 ??Precision over approximation.}
A key weakness of purely statistical approaches to code generation is their
reliance on learned heuristics that cannot provide correctness guarantees.
HSP-Agile addresses this by using a formal SMT solver as the conformance
oracle: test inputs are not drawn from a fixed hand-crafted bank but are
synthesised on-the-fly from the symbolic specification, and the reference
output is computed from the mathematical definition of the oracle $g_t$, not
from a human-annotated expected value.  This means that every acceptance
decision is backed by a formal semantics rather than a distributional
approximation.

\paragraph{P2 ??Bounded budget with guaranteed termination.}
Autonomous agents that do not terminate ??or that have no formal bound on the
number of LLM calls they make ??are impractical in production settings.
HSP-Agile enforces a hard budget $K$ on the number of synthesis-check
iterations.  Regardless of the difficulty of the task, the pipeline will
complete within $K$ LLM calls and $K$ formal checking passes.  In our
experiments we set $K = 5$; the pipeline never runs indefinitely.

\paragraph{P3 ??Counterexample specificity.}
When a candidate fails conformance checking, the feedback returned to the
LLM must be informative enough to guide a targeted repair rather than a
random re-generation.  Generic feedback such as ``your code is wrong'' is
insufficient.  HSP-Agile constructs feedback that includes: the concrete
failing input $x^*$, the oracle output $g_t(x^*)$, the candidate's
actual output $f_c(x^*)$, and the index of the FSF scenario that should have
been matched.  This specificity transforms the repair step into a targeted
correction problem rather than a blind regeneration.

\paragraph{P4 ??Conservative acceptance.}
Acceptance is conjunctive: a candidate is accepted only when it satisfies
\emph{both} the formal conformance criterion \emph{and} the requirement
pattern criterion.  A disjunctive predicate would allow formally correct code
that nonetheless violates a safety-critical requirement pattern (e.g.\
unconditionally returning success without inspecting the inputs) to be
accepted.  The conjunctive design ensures that no single stage acts as a
sufficient acceptance criterion in isolation.

\paragraph{P5 ??Modularity and ablation transparency.}
Because the pipeline is composed of four independently configurable stages,
each stage can be enabled or disabled via a boolean configuration flag.  This
design choice has two consequences: (i)~it allows clean \emph{ablation
studies} in which individual contributions are isolated and measured, and
(ii)~it allows operators to trade strictness for throughput in settings
where formal checking is too expensive.  Chapter~\ref{ch:experiments} exploits
this modularity to isolate the contributions of Stages 2, 3, and 4
individually.

% ---------------------------------------------------------------------------
\section{Pipeline Overview}
\label{sec:method:overview}
% ---------------------------------------------------------------------------

The HSP-Agile pipeline processes one Agile-SOFL task at a time.  The
high-level data flow is:
\begin{equation}
  \mathcal{S}_t
  \;\xrightarrow{\text{Stage 1}}\;
  c_k
  \;\xrightarrow{\text{Stage 2}}\;
  (\mathrm{Conf}(c_k,t),\;\mathit{cexs}_k)
  \;\xrightarrow{\text{Stage 3}}\;
  \mathit{feedback}_k
  \;\xrightarrow{\text{Stage 4}}\;
  \mathrm{Accept}(c_k,t),
  \label{eq:pipeline-flow}
\end{equation}
where $k \in \{1,\ldots,K\}$ indexes the current iteration of the repair
loop, $c_k$ is the $k$-th candidate implementation, $\mathrm{Conf}(c_k,t) \in
[0,1]$ is the formal conformance score, $\mathit{cexs}_k$ is the set of
counterexamples produced by Stage 2, $\mathit{feedback}_k$ is the structured
natural-language feedback assembled by Stage 3, and $\mathrm{Accept}(c_k,t)
\in \{0,1\}$ is the final binary acceptance decision.

The pipeline iterates equation~\eqref{eq:pipeline-flow} until either
$\mathrm{Accept}(c_k, t) = 1$ (early exit on success) or the budget $k = K$
is exhausted, at which point the best-effort candidate is returned with status
\textsc{Best-Effort}.

Figure~\ref{fig:pipeline-arch} summarises the four-stage architecture and
repair loop; Figure~\ref{fig:pipeline-seq} shows the message-level sequence
for one run.

\begin{figure}[t]
  \centering
  \tikzinput{diagrams/tikz/pipeline_architecture}
  \caption{HSP-Agile pipeline architecture.  Shadowed boxes denote
           processing stages; the repair loop (left arc) enables
           counterexample-guided refinement within budget $K$.}
  \label{fig:pipeline-arch}
\end{figure}

\begin{figure}[t]
  \centering
  \IfFileExists{diagrams/puml/rendered/pipeline_sequence.png}{%
    \pumlfig{diagrams/puml/rendered/pipeline_sequence}%
  }{%
    \fbox{\parbox{0.80\textwidth}{\centering\vspace{1.2cm}
      \textit{[Figure pending render: pipeline\_sequence.pdf]}\\
      Sequence diagram of one HSP-Agile pipeline run.\\
      The repair loop iterates until
      $\mathrm{Accept}(c_k,t)=1$ or budget $K$ is exhausted.
      \vspace{1.2cm}}}%
  }
  \caption{Sequence diagram of one HSP-Agile pipeline run.  The
           repair loop iterates until $\mathrm{Accept}(c_k,t) = 1$ or
           budget $K$ is exhausted.}
  \label{fig:pipeline-seq}
\end{figure}

\subsection*{Configuration parameters}

The pipeline is controlled by five configuration parameters, summarised in
Table~\ref{tab:config}.

\begin{table}[H]
  \caption{HSP-Agile pipeline configuration parameters.}
  \label{tab:config}
  \centering
  \thesistable{%
    \begin{tabularx}{\linewidth}{@{}llX@{}}
      \toprule
      Parameter & Type & Description \\
      \midrule
      \code{model}           & string  & LLM endpoint identifier \\
      $K$                    & integer & maximum repair iterations ($K \ge 1$) \\
      \code{enable\_formal}  & boolean & enable Stage 2 (SMT conformance checking) \\
      \code{enable\_patterns}& boolean & enable Stage 4 (requirement pattern guard) \\
      \code{enable\_repair}  & boolean & enable Stage 3 (counterexample feedback) \\
      \bottomrule
    \end{tabularx}%
  }
\end{table}

In the full HSP-Agile configuration all three boolean flags are \textit{true}.
Ablation variants set one or more flags to \textit{false}; this is the
mechanism used in the experimental evaluation of
Chapter~\ref{ch:experiments}.

% ---------------------------------------------------------------------------
\section{Stage 1: LLM-Based Candidate Synthesis}
\label{sec:method:synthesis}
% ---------------------------------------------------------------------------

Stage 1 maps a task specification ??possibly augmented with feedback from a
previous iteration ??to a candidate Python implementation.  It is the only
stage that involves a large language model call; all other stages are
deterministic.

\subsection{Prompt Structure}
\label{subsec:method:prompt}

The synthesis prompt is a two-part instruction following the OpenAI chat
completion format \citep{openai2023gpt4, chen2021codex}.
Its structure is parameterised by the task $t$ and the
feedback $\mathit{feedback}_{k-1}$ from the previous iteration
($\emptyset$ when $k = 1$).

\begin{definition}[Synthesis Prompt]
\label{def:prompt}
Let $t$ be an Agile-SOFL task and $k \ge 1$ a repair iteration index.
The \emph{synthesis prompt} $\mathcal{P}(t, k)$ is a two-message chat
transcript following the OpenAI completion format: a fixed \emph{system}
role message and a \emph{user} role message that carries the specification,
signature, implementation instructions, and (when $k > 1$) structured repair
feedback from the previous iteration.
Figure~\ref{fig:synthesis-prompt-chat} shows the layout; formally,
\[
  \mathcal{P}(t,k) = \langle \mathit{SystemMsg},\; \mathit{UserMsg}(t,k) \rangle
\]
where $\mathit{UserMsg}(t,1)$ contains $(\mathcal{S}_t,\, \mathit{sig}(t))$ and
$\mathit{UserMsg}(t,k)$ for $k > 1$ additionally includes
$(\mathit{cexs}_{k-1},\, \mathit{hits}_{k-1})$.
\end{definition}

\begin{figure}[t]
  \centering
  \tikzinput{diagrams/tikz/synthesis_prompt_chat}
  \caption{Synthesis prompt as a chat-style transcript.
           The system message (left) enforces top-to-bottom guard evaluation;
           user messages (right) inject the FSF specification, signature, and
           implementation instructions; when $k > 1$, a second user turn appends
           counterexamples and pattern violations from the prior iteration
           (orange accent).}
  \label{fig:synthesis-prompt-chat}
\end{figure}

Three components of the prompt design deserve explicit attention
(Figure~\ref{fig:prompt-design-cards}).

\begin{figure}[t]
  \centering
  \tikzinput{diagrams/tikz/prompt_design_cards}
  \caption{Three prompt-design mechanisms in Stage~1 synthesis~\citep{li2022agilesofl}.
           Scenario ordering reduces guard reordering; signature
           injection aligns variable names with the FSF specification;
           feedback context supplies repair guidance when $k > 1$.}
  \label{fig:prompt-design-cards}
\end{figure}

\subsection{Model Configuration}
\label{subsec:method:model}

We use the Qwen2.5-27B-Instruct model (\textsf{Qwen-27B}), accessed through an
OpenAI-compatible LLM gateway.  The model is queried with temperature $\tau = 0.2$ to
favour deterministic, specification-faithful completions over creative
variation.  A maximum token budget of 2048 output tokens is imposed; in
practice, all benchmark tasks require fewer than 1200 tokens for the function
body.

\subsection{Output Parsing}
\label{subsec:method:parsing}

The model's response is expected to contain exactly one fenced Python code
block delimited by \code{```python} and \code{```}.  The parser extracts the
block, strips the delimiters, and compiles the result using Python's built-in
\code{compile()} function.  If compilation fails (syntax error), or if no
code block is found, the candidate $c_k$ is set to a sentinel
$\bot$ that returns \code{None} on any call; this sentinel always fails Stage 2
with conformance score 0.

\subsection{Formal Definition}
\label{subsec:method:synth-formal}

We formalise Stage 1 as a stochastic function:
\begin{equation}
  c_k \;=\; \mathrm{LLM}\!\left(t,\;\mathit{feedback}_{k-1}\right),
  \label{eq:synthesis}
\end{equation}
where $\mathrm{LLM}$ denotes the composition of prompt construction, API
call, and output parsing.  Because the model is queried with $\tau = 0.2$
(non-zero temperature), $c_k$ is a random variable; the formal checker in
Stage 2 provides the deterministic conformance verdict that makes the
pipeline's acceptance decision temperature-independent.

% ---------------------------------------------------------------------------
\section{Stage 2: SMT-Driven Formal Conformance Checking}
\label{sec:method:checker}
% ---------------------------------------------------------------------------

Stage 2 is the most novel contribution of HSP-Agile \citep{demoura2008z3, barrett2018smt}.
Given candidate $c_k$ and task $t$, it computes $\mathrm{Conf}(c_k, t) \in
[0,1]$ and counterexamples $\mathit{cexs}_k$ explaining every failure.
The stage consists of four steps, summarised in Algorithm~\ref{alg:checker}.

\subsection{Step 1: Scenario-to-Constraint Translation}
\label{subsec:method:translate}

For each scenario $s_i = (g_i, \phi_i)$ in $\mathcal{S}_t$, the guard $g_i$
is first parsed by a recursive descent parser that recognises linear arithmetic
expressions over integer variables and their Boolean combinations.  The parser
produces an abstract syntax tree (AST) that is then translated into a Z3
constraint.

\begin{definition}[Z3 Guard Encoding]
\label{def:z3-guard}
Let $\mathcal{X}_t = \mathbb{Z}^d$ with variable names $v_1,\ldots,v_d$.
For each $v_j$ we create a Z3 integer constant $\hat{v}_j$.  The
\emph{Z3 encoding} of guard $g_i$ is a formula $\hat{g}_i$ obtained by
structural recursion:
\begin{align*}
  \widehat{a > b}   &\;\triangleq\; \hat{a} > \hat{b} \\
  \widehat{a \ge b} &\;\triangleq\; \hat{a} \ge \hat{b} \\
  \widehat{a = b}   &\;\triangleq\; \hat{a} = \hat{b} \\
  \widehat{p \land q} &\;\triangleq\; \mathtt{And}(\hat{p}, \hat{q}) \\
  \widehat{p \lor q}  &\;\triangleq\; \mathtt{Or}(\hat{p}, \hat{q}) \\
  \widehat{\lnot p}   &\;\triangleq\; \mathtt{Not}(\hat{p}).
\end{align*}
\end{definition}

To respect first-match semantics, the constraint for scenario $s_i$ is not
simply $\hat{g}_i$ but the \emph{priority-constrained formula}:
\begin{equation}
  \Phi_i \;\triangleq\; \hat{g}_i \;\land\;
    \bigwedge_{j=1}^{i-1} \lnot\,\hat{g}_j,
  \label{eq:priority-constraint}
\end{equation}
which ensures that the satisfying model generated for scenario $i$ does not
also satisfy any earlier guard ??i.e.\ the concrete input genuinely activates
scenario $i$ under first-match evaluation.

\subsection{Step 2: Satisfying Model Generation}
\label{subsec:method:z3-solve}

For each priority-constrained formula $\Phi_i$, an independent Z3 solver
instance is created and queried:

\inputpython{listings/z3_gen.py}{Z3 test-case generation for scenario~$i$.}{lst:z3-gen}

If $\Phi_i$ is unsatisfiable ??which can happen when the guard combination is
logically contradictory ??the scenario contributes no test case to
$\mathcal{D}(t)$.  In practice this occurs only for the ``others'' scenario
when all earlier guards collectively cover the entire input domain; the
solver returns \textsf{unsat} and the scenario is skipped.

\subsection{Step 3: Oracle Output Computation}
\label{subsec:method:oracle-compute}

For each concrete input $x_i \in \mathcal{D}(t)$, the oracle output
$y_i^* = g_t(x_i)$ is computed by evaluating the FSF output mapping
$\phi_i$ symbolically.  Because output mappings in Agile-SOFL are conjunctions
of assignment statements (e.g.\ $\mathit{risk} := \texttt{'HIGH'},\;
\mathit{action} := \texttt{'ESCALATE'},\; \mathit{score} := 90$), the
evaluation is straightforward substitution of the concrete values of
$x_i$ into the right-hand sides.  The result is a Python dictionary
$\{$\textit{field}: \textit{value}$\}$ that serves as the reference output.

\subsection{Step 4: Candidate Execution and Comparison}
\label{subsec:method:exec}

The candidate $c_k$ is executed in a restricted sandbox: the function body is
compiled with \code{exec()} into an isolated namespace, the concrete input
$x_i$ is passed as keyword arguments, and the returned value is compared
field-by-field with the oracle output $y_i^*$.

Comparison is structural: for each output field $f$, we check
$\hat{y}[f] = y^*[f]$ using Python's equality operator after type
normalisation (string case normalisation and numeric rounding to two decimal
places).  A candidate passes scenario $i$ if and only if all fields match.

\begin{definition}[Conformance Score]
\label{def:conf-score}
Given candidate $c$ and task $t$ with case set $\mathcal{D}(t)$, the
\emph{conformance score} is:
\begin{equation}
  \mathrm{Conf}(c, t) \;=\;
    \frac{\bigl|\{x \in \mathcal{D}(t) : f_c(x) = g_t(x)\}\bigr|}
         {|\mathcal{D}(t)|},
  \label{eq:conf-score}
\end{equation}
where $f_c(x)$ denotes the output of candidate $c$ on input $x$.
$\mathrm{Conf}(c, t) = 1$ if and only if $c$ passes every generated test case.
\end{definition}

\subsection{Algorithm}
\label{subsec:method:checker-alg}

Algorithm~\ref{alg:checker} presents the complete formal conformance checker.

\begin{algorithm}[t]
\caption{SMT-Driven Formal Conformance Checking}
\label{alg:checker}
\begin{algorithmic}[1]
\Require task $t = (\mathcal{S}_t, \mathcal{X}_t, \mathcal{Y}_t)$, candidate~$c$
\Ensure conformance score in $[0,1]$, counterexample list $\mathit{cexs}$
\State $\mathcal{D}(t) \leftarrow \emptyset$;\quad $\mathit{cexs} \leftarrow \emptyset$
\For{$i = 1$ \textbf{to} $n$}
  \State $\Phi_i \leftarrow \ProcCall{TranslateToZ3}{s_i,\, \{g_j\}_{j<i}}$
         \Comment{eq.~\eqref{eq:priority-constraint}}
  \If{$\ProcCall{Z3-Solve}{\Phi_i} = \mathit{sat}$}
    \State $x_i \leftarrow \ProcCall{ExtractModel}{}$
    \State $y_i^* \leftarrow \phi_i(x_i)$
    \State $\mathcal{D}(t) \leftarrow \mathcal{D}(t) \cup \{(x_i, y_i^*)\}$
  \EndIf
\EndFor
\State $\mathit{passed} \leftarrow 0$
\For{\textbf{each} $(x, y^*) \in \mathcal{D}(t)$}
  \State $y^c \leftarrow f_c(x)$
  \If{$y^c = y^*$}
    \State $\mathit{passed} \leftarrow \mathit{passed} + 1$
  \Else
    \State $\mathit{cexs} \leftarrow \mathit{cexs} \cup
           \{(x, y^*, y^c, \ProcCall{ScenarioOf}{x})\}$
  \EndIf
\EndFor
\State \Return $(\mathit{passed}\,/\,|\mathcal{D}(t)|,\,\mathit{cexs})$
\end{algorithmic}
\end{algorithm}

\subsection{Why SMT Over Random Testing}
\label{subsec:method:smt-rationale}

The choice to use an SMT solver rather than random or boundary-value testing
follows from the low probability of hitting multi-conjunct guard regions by
chance \citep{claessen2000quickcheck, zhu1997specification}.
The priority-constrained formula $\Phi_i$ (equation~\eqref{eq:priority-constraint})
guarantees each test input lies in scenario $i$'s exclusive region under
first-match semantics, eliminating oracle ambiguity.

% ---------------------------------------------------------------------------
\section{Stage 3: Counterexample-Guided Repair}
\label{sec:method:repair}
% ---------------------------------------------------------------------------

When Stage 2 finds that $\mathrm{Conf}(c_k, t) < 1$, or when Stage 4 finds
pattern violations, Stage 3 assembles structured feedback for the next
synthesis call \citep{lynn2024verifierloop, solarlezama2008sketching}.
Repair quality is bounded by feedback specificity: the prompt must identify
\emph{which scenario} governs the failing case and supply the witness input.

\subsection{Feedback Data Structure}
\label{subsec:method:feedback-struct}

\begin{definition}[Feedback Record]
\label{def:feedback}
A \emph{feedback record} $\mathit{fb}_k$ at iteration $k$ is a structured
object:
\[
  \mathit{fb}_k \;=\;
  \left\{
  \begin{aligned}
    \mathit{counterexamples}     &\mapsto \mathit{cexs}_k,\\
    \mathit{pattern\_violations} &\mapsto \mathit{hits}_k,\\
    \mathit{attempt}            &\mapsto k,\\
    \mathit{conformance}        &\mapsto \mathrm{Conf}(c_k, t)
  \end{aligned}
  \right\}.
\]
Each counterexample in $\mathit{cexs}_k$ is a tuple
$(x^*, y^*, y^c, i^*)$ where $x^*$ is the failing input, $y^*$ is the
expected oracle output, $y^c$ is the candidate's actual output, and $i^*$ is
the scenario index whose priority-constrained region contains $x^*$.
Each pattern hit in $\mathit{hits}_k$ is a record $\{$\textit{pattern\_id},
\textit{line}, \textit{severity}, \textit{message}$\}$.
\end{definition}

\subsection{Feedback Verbalisation}
\label{subsec:method:verbalise}

Raw tuples are not maximally useful to the LLM.  Stage 3 verbalises each
counterexample into a sentence of the form:

\begin{quote}
  \textit{``For inputs $(a=15,\; b=8,\; c=3,\; d=0,\; e=0)$, your code
  returned \{\textit{risk}=\code{MEDIUM}, \textit{action}=\code{REVIEW},
  \textit{score}=60\}, but the specification requires
  \{\textit{risk}=\code{HIGH}, \textit{action}=\code{ESCALATE},
  \textit{score}=90\}.
  This input satisfies Scenario S1 ($a > 10$ AND $b > 5$), which should be matched first.''}
\end{quote}

This sentence encodes all four components of the counterexample tuple:
the input values, the expected and actual outputs side-by-side, and the
causal scenario index with its guard condition quoted verbatim from the FSF.
Controlled experiments (not reported in this thesis) show that including the
scenario guard in the feedback text significantly reduces the number of
repair iterations needed.

Pattern violation feedback is similarly verbalised.

Figure~\ref{fig:repair-feedback} summarises how counterexamples and pattern
hits are packaged into the next synthesis prompt.
Figure~\ref{fig:cegis-loop} situates this iteration within the classical
CEGIS synthesiser--verifier loop.

\begin{figure}[t]
  \centering
  \tikzinput{diagrams/tikz/repair_feedback}
  \caption{Counterexample-guided repair iteration.  Each failed attempt
           produces structured feedback combining formal counterexamples
           and pattern violations, used to guide the next synthesis
           attempt.}
  \label{fig:repair-feedback}
\end{figure}

\begin{figure}[t]
  \centering
  \tikzinput{diagrams/tikz/cegis_loop}
  \caption{The CEGIS feedback loop as instantiated in HSP-Agile.  The LLM
           serves as the synthesiser, the formal checker as the verifier,
           and structured feedback bridges the two.}
  \label{fig:cegis-loop}
\end{figure}

\subsection{Relation to CEGIS}
\label{subsec:method:cegis}

The repair loop is an instance of the \emph{counterexample-guided inductive
synthesis} (CEGIS) paradigm \citep{solarlezama2008sketching}, adapted for the
LLM-as-synthesiser setting.  In classical CEGIS, the synthesiser enumerates
programs from a grammar and the verifier is a model checker; in HSP-Agile,
the synthesiser is a large language model and the verifier is a Z3-backed
conformance checker.  The structural correspondence is exact:
\begin{itemize}
  \item \textbf{Synthesiser}: $\mathrm{LLM}(t, \mathit{fb}_{k-1})$
        produces a candidate $c_k$ conditioned on prior feedback.
  \item \textbf{Verifier}: $\mathrm{FormalCheck}(c_k, t)$ either accepts
        $c_k$ (conformance 1) or produces a counterexample.
  \item \textbf{Feedback bridge}: Stage 3 translates the formal
        counterexample into natural language that the LLM can interpret.
\end{itemize}

The key difference from classical CEGIS is that the synthesiser is not
complete: the LLM may fail to find a correct implementation even given a
perfect counterexample, because it does not enumerate a finite grammar.
The bounded budget $K$ acknowledges this limitation and ensures practical
termination.

% ---------------------------------------------------------------------------
\section{Stage 4: Requirement Pattern Guard}
\label{sec:method:guard}
% ---------------------------------------------------------------------------

Even a candidate with $\mathrm{Conf}(c, t) = 1$ may violate structural
requirements not captured by the finite test suite $\mathcal{D}(t)$ \citep{hovemeyer2004findbugs, li2023iet}.
Stage~4 guards against such pathologies using a catalogue of requirement error
patterns; Table~\ref{tab:patterns} lists all 14 patterns and their severities.

\subsection{Pattern Catalog}
\label{subsec:method:catalog}

The HSP-Agile pattern catalog $\mathcal{P}$ contains 14 patterns formalising
implementation faults observed in LLM outputs on Agile-SOFL tasks.

\begin{table}[H]
  \caption{HSP-Agile Error Pattern Catalog ($|\mathcal{P}| = 14$).
           Severity follows the convention: \textbf{critical} patterns
           cause immediate acceptance failure; \textit{high} patterns
           are included in the high-severity set~$\mathcal{P}^H$.}
  \label{tab:patterns}
  \centering
  \thesistable{%
    \small
    \setlength{\tabcolsep}{4pt}
    \begin{tabularx}{\linewidth}{@{}lXll@{}}
      \toprule
      ID   & Name                     & Category       & Severity \\
      \midrule
      RF01 & Missing precondition check  & requirement    & high \\
      RF02 & Wrong default branch        & FSF            & high \\
      RF03 & Hardcoded magic number      & implementation & medium \\
      RF04 & Missing output field        & interface      & high \\
      RF05 & Swapped comparison          & logic          & high \\
      RF06 & No guard for zero input     & boundary       & high \\
      RF07 & Unconditional success       & logic          & \textbf{critical} \\
      RF08 & Missing type coercion       & implementation & medium \\
      RF09 & Incomplete FSF coverage     & FSF            & high \\
      RF10 & External variable ignored   & SOFL           & medium \\
      RF11 & Off-by-one boundary         & boundary       & high \\
      RF12 & Duplicate assignment        & implementation & low \\
      RF13 & Empty handler               & logic          & high \\
      RF14 & Spec-comment mismatch       & requirement    & medium \\
      \bottomrule
    \end{tabularx}%
  }
\end{table}

\paragraph{Interface patterns (RF04, RF10).}
\emph{RF04} (Missing output field) fires when the function's return value
omits a field that is declared in the FSF output domain $\mathcal{Y}_t$.
\emph{RF10} (External variable ignored) fires when the specification
references an external program variable (a SOFL notion) but the implementation
does not read or modify it.

\paragraph{Logic patterns (RF05, RF07, RF13).}
\emph{RF05} (Swapped comparison) fires when a comparison operator is inverted
relative to the specification ??e.g.\ $a < 10$ in the code where $a > 10$
is required.  Detection combines guard extraction from the FSF with AST
analysis of the candidate.
\emph{RF07} (Unconditional success) is the most critical pattern: it fires
when the function returns a fixed value on any code path that reaches the
return statement without evaluating at least one guard-dependent branch.
\emph{RF13} (Empty handler) fires when a branch corresponding to an FSF
scenario contains only a \code{pass} statement or an empty body.

\paragraph{FSF conformance patterns (RF02, RF09).}
\emph{RF02} (Wrong default branch) fires when the implementation's
``otherwise'' branch returns values inconsistent with the final FSF scenario
$s_n$.  \emph{RF09} (Incomplete FSF coverage) fires when the number of
distinct conditional branches in the implementation is strictly less than $n-1$
(the number of non-default scenarios).

\paragraph{Boundary patterns (RF06, RF11).}
\emph{RF06} (No guard for zero input) fires when a numeric guard in the FSF
includes zero as a boundary value but the implementation does not handle the
$= 0$ case explicitly.  \emph{RF11} (Off-by-one boundary) fires when
a threshold comparison uses $\ge$ where $>$ is specified or vice versa,
detectable by comparing extracted threshold literals with FSF guard constants.

\paragraph{Requirement patterns (RF01, RF14).}
\emph{RF01} (Missing precondition check) fires when the FSF specification
includes an explicit precondition block but the implementation does not include
corresponding guard logic.  \emph{RF14} (Spec-comment mismatch) fires when
a function docstring or inline comment references a scenario condition that
differs from the actual code logic, indicating a copy-paste error.

\paragraph{Implementation patterns (RF03, RF08, RF12).}
These three medium-to-low severity patterns detect coding-quality issues:
magic numbers embedded in the function body rather than derived from the
specification, missing type casts when the spec requires integer output but
the implementation returns floats, and duplicate variable assignments that
shadow earlier values.

\subsection{Detection Method}
\label{subsec:method:detection}

Pattern detection uses a hybrid of Python AST analysis and regular expression
matching over the source text.  For each pattern $p \in \mathcal{P}$, a
detector function $\delta_p : \mathcal{C} \times \mathcal{T} \to \{0,1\}$
takes the candidate source $c$ and task $t$ as inputs and returns 1 if the
pattern fires.  The detectors for structural patterns (RF04, RF07, RF09,
RF13) operate on the AST; those for surface patterns (RF03, RF08, RF14)
use regex over source text; and those for guard-comparison patterns
(RF01, RF02, RF05, RF06, RF11) cross-reference the extracted FSF guards with
the candidate's conditional structure.

\begin{definition}[Pattern Hit Set]
\label{def:pattern-hits}
Given candidate $c$ and task $t$, the \emph{pattern hit set} is:
\[
  \mathrm{Hits}(c, t) \;=\;
    \bigl\{p \in \mathcal{P} : \delta_p(c, t) = 1\bigr\}.
\]
\end{definition}

\subsection{High-Severity Subset}
\label{subsec:method:high-sev}

Not all patterns carry equal weight in the acceptance decision.  We define
the \emph{high-severity set}:
\begin{equation}
  \mathcal{P}^H \;=\;
    \{\text{RF01},\;\text{RF02},\;\text{RF04},\;\text{RF05},\;\text{RF06},\;
      \text{RF07},\;\text{RF09},\;\text{RF11},\;\text{RF13}\},
  \label{eq:high-sev}
\end{equation}
comprising the nine patterns whose severity is rated \emph{high} or
\emph{critical} in Table~\ref{tab:patterns}.  A candidate is not accepted
if any $p \in \mathcal{P}^H$ fires, regardless of its conformance score.
Medium- and low-severity patterns are recorded and reported but do not
block acceptance; they inform the feedback message for the next iteration.

% ---------------------------------------------------------------------------
\section{Hybrid Acceptance Predicate}
\label{sec:method:accept}
% ---------------------------------------------------------------------------

The outputs of Stages 2 and 4 are combined into a single binary acceptance
decision by the \emph{hybrid acceptance predicate}.

\begin{definition}[Hybrid Acceptance Predicate]
\label{def:accept}
Given candidate $c$ and task $t$, the \emph{hybrid acceptance predicate} is:
\begin{equation}
  \mathrm{Accept}(c, t) \;=\;
    \underbrace{\mathbf{1}\!\left[\mathrm{Conf}(c, t) = 1\right]}_{\text{formal correctness}}
    \;\land\;
    \underbrace{\mathbf{1}\!\left[\sum_{p \in \mathcal{P}^H}
      \delta_p(c, t) = 0\right]}_{\text{requirement safety}}.
  \label{eq:accept}
\end{equation}
\end{definition}

The predicate is conjunctive: both conditions must hold simultaneously.
This design choice reflects Principle~P4 (Section~\ref{sec:method:principles})
and is justified by two complementary observations.

\paragraph{Why formal conformance alone is insufficient.}
The formal checker generates exactly $n$ test cases ??one per FSF scenario ??from which each covers only a single point in the scenario's exclusive guard
region.  A pathological implementation that hard-codes the correct output for
each of these $n$ points would achieve $\mathrm{Conf}(c,t) = 1$ while being
completely wrong on all other inputs.  Stage 4 detects this via pattern RF07
(Unconditional success) and RF03 (Hardcoded magic number).

\paragraph{Why pattern guard alone is insufficient.}
The pattern catalog is necessarily finite and incomplete: it catches known
anti-patterns but cannot guarantee conformance for scenarios it does not
model.  A candidate that passes all 14 pattern detectors but implements the
guard conditions in the wrong order would satisfy the pattern guard but fail
the formal conformance check on inputs where two guards overlap.

\paragraph{Complementarity.}
The two conditions are therefore genuinely complementary: formal conformance
certifies correctness on the generated test suite; the pattern guard certifies
structural soundness with respect to known fault classes.  Together they form
a conservative but not overly restrictive acceptance criterion.

\subsection{Full Pipeline Algorithm}
\label{subsec:method:full-alg}

Algorithm~\ref{alg:pipeline} presents the complete HSP-Agile pipeline,
incorporating all four stages and the acceptance predicate.

\begin{algorithm}[t]
\caption{HSP-Agile Candidate Synthesis and Fault Prevention}
\label{alg:pipeline}
\begin{algorithmic}[1]
\Require task $t$; config $=$ ($\mathit{model}$, $K$, $\mathit{enable\_formal}$,
         $\mathit{enable\_patterns}$, $\mathit{enable\_repair}$)
\Ensure accepted candidate $c^*$ and status $\in \{\textsc{Accept},
        \textsc{Best-Effort}\}$
\State $\mathit{feedback} \leftarrow \emptyset$;\quad
       $\mathit{best\_conf} \leftarrow 0$;\quad
       $c^* \leftarrow \bot$
\For{$k = 1$ \textbf{to} $K$}
  \State $c_k \leftarrow \ProcCall{LLM}{t,\, \mathit{feedback}}$
         \Comment{Stage 1 (eq.~\eqref{eq:synthesis})}
  \If{$\mathit{enable\_formal}$}
    \State $\mathit{conf},\;\mathit{cexs} \leftarrow
           \ProcCall{FormalCheck}{c_k,\, t}$
           \Comment{Stage 2 (Algorithm~\ref{alg:checker})}
  \Else
    \State $\mathit{conf} \leftarrow \ProcCall{UnitTestScore}{c_k,\, t}$;\quad
           $\mathit{cexs} \leftarrow \emptyset$
  \EndIf
  \If{$\mathit{conf} > \mathit{best\_conf}$}
    \State $\mathit{best\_conf} \leftarrow \mathit{conf}$;\quad
           $c^* \leftarrow c_k$
  \EndIf
  \If{$\mathit{conf} = 1$}
    \If{$\mathit{enable\_patterns}$}
      \State $\mathit{hits} \leftarrow \ProcCall{PatternGuard}{c_k,\, t}$
             \Comment{Stage 4}
      \State $\mathit{hits}^H \leftarrow
             \{h \in \mathit{hits} : h.\mathit{severity} \ge \mathit{high}\}$
      \If{$|\mathit{hits}^H| = 0$}
        \State \Return $c_k$,\enspace \textsc{Accept}
               \Comment{predicate~\eqref{eq:accept} satisfied}
      \EndIf
    \Else
      \State \Return $c_k$,\enspace \textsc{Accept}
    \EndIf
  \Else
    \State $\mathit{hits} \leftarrow \emptyset$
  \EndIf
  \If{$\mathit{enable\_repair}$}
    \State $\mathit{feedback} \leftarrow
           \ProcCall{BuildFeedback}{\mathit{cexs},\,\mathit{hits},\,k}$
           \Comment{Stage 3}
  \EndIf
\EndFor
\State \Return $c^*$,\enspace \textsc{Best-Effort}
\end{algorithmic}
\end{algorithm}

\paragraph{Reading the algorithm.}
Lines 1--3 initialise the loop state.  Lines 4--9 call Stage 1 and Stage 2
(or a simpler unit-test scorer if \textit{enable\_formal} is false).  Lines
10--11 track the best candidate seen so far by conformance score, so that the
best-effort fallback (line 23) returns the most promising candidate if no
accepted solution is found.  Lines 12--21 implement the conjunctive acceptance
check: Stage 4 is only invoked when conformance is perfect ($\mathit{conf} =
1$), avoiding the cost of pattern analysis on definitively failing candidates.
Line 22 builds the feedback record for the next iteration.

% ---------------------------------------------------------------------------
\section{Theoretical Properties}
\label{sec:method:theory}
% ---------------------------------------------------------------------------

We now establish the formal properties of the HSP-Agile pipeline.  The proofs
are elementary given the algorithm structure; their value is in making the
guarantees explicit and verifiable.

\begin{theorem}[Termination]
\label{thm:termination}
Algorithm~\ref{alg:pipeline} terminates in at most $K$ iterations,
making at most $K$ LLM API calls and at most $K$ invocations of
Algorithm~\ref{alg:checker}.
\end{theorem}

\begin{proof}
The body of the \textbf{for} loop at line 4 either returns at line 16
or 18 (\textsc{Accept}) or completes and increments the loop counter.
The loop counter is bounded above by $K$ by the loop guard.  No recursive
calls or unbounded sub-loops appear in either algorithm.  Therefore
Algorithm~\ref{alg:pipeline} terminates after at most $K$ complete
iterations.  Similarly, Algorithm~\ref{alg:checker} iterates over the
finite sequence $\mathcal{S}_t$ (length $n$) and over the finite set
$\mathcal{D}(t)$ (size $|\mathcal{D}(t)| \le n$); both loops terminate.
\end{proof}

\begin{theorem}[Soundness of Acceptance]
\label{thm:soundness}
If Algorithm~\ref{alg:pipeline} returns $c_k$ with status \textsc{Accept},
then $\mathrm{Conf}(c_k, t) = 1$ and
$|\{p \in \mathcal{P}^H : \delta_p(c_k, t) = 1\}| = 0$.
\end{theorem}

\begin{proof}
Status \textsc{Accept} is returned at line 16 or 18.  Line 16 is reached
only when the condition at line 13 ($\mathit{conf} = 1$) and the condition
at line 15 ($|\mathit{hits}^H| = 0$) are both satisfied.  By
Definition~\ref{def:conf-score}, $\mathit{conf} = 1$ is equivalent to
$\mathrm{Conf}(c_k, t) = 1$.  By Definition~\ref{def:pattern-hits} and
equation~\eqref{eq:high-sev}, $|\mathit{hits}^H| = 0$ is equivalent to
$|\{p \in \mathcal{P}^H : \delta_p(c_k,t)=1\}| = 0$.  Line 18 is reached
when \textit{enable\_patterns} is false; in that case the pattern condition
is vacuously satisfied.  The theorem follows from the conjunction of these
two conditions, which is exactly the acceptance predicate~\eqref{eq:accept}.
\end{proof}

\begin{proposition}[Case Coverage]
\label{prop:coverage}
Under the assumption that $\textsc{Z3-Solve}(\Phi_i)$ returns \textsf{sat}
for every $i \in \{1, \ldots, n\}$, the case set generated by
Algorithm~\ref{alg:checker} satisfies $|\mathcal{D}(t)| = n$.
\end{proposition}

\begin{proof}
Under the stated assumption, each of the $n$ iterations of the first loop
in Algorithm~\ref{alg:checker} appends exactly one element to $\mathcal{D}(t)$.
The elements are distinct because the priority-constrained formula $\Phi_i$
(equation~\eqref{eq:priority-constraint}) excludes the guard regions of all
earlier scenarios, so the satisfying models for different $i$ lie in disjoint
regions of $\mathcal{X}_t$.
\end{proof}

\begin{remark}
In practice the assumption of Proposition~\ref{prop:coverage} can fail for
the last scenario $s_n$ when the earlier guards collectively cover
$\mathcal{X}_t$.  This occurs in the benchmark when all integer-valued inputs
satisfy one of the first $n-1$ guards; the solver returns \textsf{unsat} for
$\Phi_n$ and the ``others'' case contributes no test case.  In that setting
$|\mathcal{D}(t)| = n - 1$, and conformance is computed over $n-1$ cases.
This does not affect the soundness guarantee of
Theorem~\ref{thm:soundness} ??it merely reduces the strength of the
test suite for that particular scenario.
\end{remark}

\subsection{Bounded Prevention Guarantee}
\label{subsec:method:prevention}

The two theorems above together define the \emph{bounded prevention guarantee}
of HSP-Agile:

\begin{corollary}[Bounded Fault Prevention]
\label{cor:prevention}
Every candidate returned with status \textsc{Accept} by
Algorithm~\ref{alg:pipeline} is free of all faults detectable by
$\mathcal{D}(t)$ and all high-severity patterns in $\mathcal{P}^H$.
The guarantee holds within $K$ LLM calls regardless of the task.
\end{corollary}

This corollary is the formal counterpart of the thesis's central claim: that
HSP-Agile \emph{prevents} rather than merely \emph{detects} specification
faults, within a bounded computational budget.  The prevention is not
universal ??it is bounded by the completeness of $\mathcal{D}(t)$ and
$\mathcal{P}^H$ ??but it is decidable, transparent, and auditable in a way
that post-hoc statistical fault detection is not.

\subsection{Computational Complexity}
\label{subsec:method:complexity}

The dominant cost per iteration is the LLM API call (Stage 1).
The formal checker (Stage 2) makes $n$ Z3 solver calls, each polynomial for
linear arithmetic \citep{barrett2018smt, demoura2008z3}.
Pattern detection (Stage 4) runs in $O(|c|)$ where $|c|$ is the candidate
source length \citep{cousot1977abstract}.
Feedback verbalisation (Stage 3) is $O(|\mathit{cexs}|) = O(n)$.
Overall pipeline complexity is $O(K \cdot n)$ formal solver calls, manageable
for benchmark sizes with $n \le 8$.

\subsection{Summary of Stage Contributions}
\label{subsec:method:stage-summary}

Table~\ref{tab:stage-summary} summarises the role, inputs, outputs, and
formal properties of each pipeline stage.

\begin{inlinetable}{Summary of HSP-Agile pipeline stages.}
  \label{tab:stage-summary}
  \thesistable{%
    \footnotesize
    \setlength{\tabcolsep}{3pt}
    \begin{tabularx}{\linewidth}{@{}clXXl@{}}
      \toprule
      Stage & Name          & Input               & Output                     & Property \\
      \midrule
      1     & Synthesis     & $t$, $\mathit{fb}$  & $c_k$                      & stochastic \\
      2     & Formal Check  & $c_k$, $t$          & $\mathrm{Conf}$, $\mathit{cexs}$ & sound (Thm.~\ref{thm:soundness}) \\
      3     & Repair        & $\mathit{cexs}$, $\mathit{hits}$ & $\mathit{fb}_k$ & CEGIS bridge \\
      4     & Pattern Guard & $c_k$, $t$          & $\mathit{hits}$            & finite catalog \\
      5     & Accept        & $\mathrm{Conf}$, $\mathit{hits}$ & $\{0,1\}$     & conjunctive \\
      \bottomrule
    \end{tabularx}%
  }
\end{inlinetable}

% ---------------------------------------------------------------------------
\section*{Chapter Summary}
% ---------------------------------------------------------------------------

This chapter has presented the complete design and formal theory of the
HSP-Agile pipeline.  Five design principles motivate a four-stage architecture:
LLM synthesis (Stage~1), SMT-driven conformance checking (Stage~2),
counterexample-guided repair (Stage~3), and requirement pattern guarding
(Stage~4).  The hybrid acceptance predicate combines Stages~2 and~4
conjunctively.

Theorems~\ref{thm:termination} and \ref{thm:soundness} establish termination
within $K$ iterations and soundness of the acceptance decision.
Corollary~\ref{cor:prevention} crystallises the bounded fault-prevention
guarantee that is the central claim of this thesis.

Chapter~\ref{ch:implementation} describes the software implementation of
this design in detail; Chapter~\ref{ch:experiments} evaluates the pipeline
experimentally on the 20-task HSP-Agile benchmark suite.
